% Generated draft: requires mathematical and source review.
\documentclass[11pt]{article}
\usepackage[margin=25mm]{geometry}
\usepackage{fontspec}
\setmainfont{Latin Modern Roman}
\usepackage{amsmath,amssymb,mathtools,unicode-math}
\setmathfont{Latin Modern Math}
\usepackage{enumitem,needspace,xurl,hyperref}
\setlength{\parindent}{0pt}
\setlength{\parskip}{5pt}
\setlength{\emergencystretch}{3em}
\newcommand{\N}{\mathbb{N}}
\newcommand{\Z}{\mathbb{Z}}
\newcommand{\Q}{\mathbb{Q}}
\newcommand{\R}{\mathbb{R}}
\newcommand{\C}{\mathbb{C}}
\newcommand{\F}{\mathbb{F}}
\newcommand{\A}{\mathbb{A}}
\newcommand{\PP}{\mathbb{P}}
\newcommand{\E}{\mathbb{E}}
\newcommand{\eps}{\varepsilon}
\newcommand{\dd}{\,\mathrm{d}}
\begin{document}
\section*{UnsolvedMath ID: 30006350}
\textbf{Source problem code:} catalogue entry
\subsection{Definitions and mathematical statement}
Let \(p\) be a prime. The source associates a finite group \(G_{E,P}(\F_p)\) of nilpotency class two to a pair \((E,P)\) involving an elliptic curve and a point. However, the supplied record does not define this construction, the permitted pairs, or the family denoted by \(\mathcal G_p\) in the original statement. These are unresolved prerequisites to a complete formulation; no domain for \(P\) or restriction on \(p\) is imposed here without source evidence.

For a group \(H\), define its lower central series by \(\gamma_1(H)=H\) and \(\gamma_{i+1}(H)=[\gamma_i(H),H]\). Thus \(H\) has nilpotency class exactly three when \(\gamma_4(H)=1\) and \(\gamma_3(H)\ne1\). A central extension in the question means an exact sequence
\[
1\longrightarrow \Z/p\Z\xrightarrow{\iota}H\xrightarrow{\pi}G_{E,P}(\F_p)\longrightarrow1,
\qquad \iota(\Z/p\Z)\subseteq Z(H).
\]
The first requested classification is: among the pairs allowed by the source construction, determine exactly those \((E,P)\) for which such an extension exists with
\[
|H|=p^{10},\qquad \gamma_4(H)=1,\qquad \gamma_3(H)\ne1.
\]
Exactness then necessarily gives \(|G_{E,P}(\F_p)|=p^9\). The source calls these extensions immediate descendants. Its precise convention for that term, and its relationship to any additional descendant conditions, are not supplied.

For an explicit interpretation of counting descendants as isomorphism classes of groups, define, for a finite group \(K\),
\[
d_p(K)=\#\left\{[H]\;\middle|\;
\begin{gathered}
|H|=p^{10},\quad \gamma_4(H)=1,\quad \gamma_3(H)\ne1,\\
\exists C\le Z(H):\ |C|=p,\quad H/C\cong K
\end{gathered}
\right\}.
\]
Here \([H]\) denotes the abstract group isomorphism class, so different choices of \(C\) in the same group do not contribute additional terms. This definition makes the extension-counting interpretation explicit; agreement with the intended immediate-descendant convention requires review.

Once the missing family \(\mathcal G_p\) is specified, the second question, under this counting interpretation, asks whether
\[
D(p)=\sum_{[K]\in\mathcal G_p/\cong}d_p(K)
\]
is polynomial on residue classes, abbreviated PORC. For a function on the admissible primes, this means that there exist an integer \(m\ge1\) and polynomials \(f_0,\ldots,f_{m-1}\in\Q[X]\) such that
\[
D(p)=f_r(p)\qquad\text{whenever }p\text{ is admissible and }p\equiv r\pmod m.
\]
The admissible primes, and whether finitely many exceptional primes may be omitted, remain unspecified. The sum counts each base-group isomorphism class once; it does not identify contributions from different base groups if they share a descendant.

As background, the original statement attributes to du Sautoy and Vaughan-Lee a non-PORC immediate-descendant count for the construction with \(E:y^2=x^3-x\) and \(P=(0,0)\). This is a supplied attribution, with its characteristic restrictions and counting conventions unavailable here. The original statement describes PORC behavior of the aggregate sum as a possibility of interest, not an established assertion.
\subsection{Short English statement}
Classify the permitted elliptic-curve pairs \((E,P)\) whose associated group has a central extension of order \(p^{10}\) and nilpotency class exactly three. Then ask whether the sum of descendant counts over the intended base-group isomorphism classes is PORC in \(p\). The construction, permitted pairs, prime restrictions, family \(\mathcal G_p\), and precise counting convention are missing, so this is a conditional statement-only draft requiring human mathematical review.
\subsection{Sources}
\begin{itemize}
\item Local UnsolvedMath registry, supplied record with permanent ID \(30006350\), title \emph{Immediate Descendants of Elliptic-Curve p-Groups}, source problem code \texttt{catalogue entry}, in the Oberwolfach Reports Open Problems collection. The record gives proposed year 2025 and identifies no individual proposer.
\item Supplied report citation: \emph{Computational Group Theory} (2025), Oberwolfach Reports, DOI \texttt{10.4171/owr/2025/27}, pp.\ 1389--1450. \url{https://doi.org/10.4171/owr/2025/27}. The supplied background describes these reports as records of workshops and their open-problem discussions.
\item Supplied literature reference: du Sautoy--Vaughan-Lee, \emph{Non-PORC descendant p-groups}. \url{https://www.sciencedirect.com/science/article/pii/S0021869312001901}. The original excerpt attributes the exceptional elliptic-curve example to these authors using reference label [1]; the supplied literature list identifies this paper.
\item Supplied literature reference: \emph{Oberwolfach Report 27/2025}. \url{https://ems.press/content/serial-article-files/51860}.
\item Literature assessment dated exactly \texttt{2026-08-22}: the supplied curation labels the problem open, reports a known exceptional family with non-PORC descendant counts, and states that a classification of pairs admitting the requested extension is not known. This is supplied dated curation derived from the OWR material, not independently verified literature assessment; the source itself requests independent verification of time-sensitive status and cited claims.
\item The registry labels the statement \texttt{corrected\_verified}, explaining that only formula delimiters, notation, grammar, and typography were repaired while mathematical content and scope were preserved. This supplied label does not establish mathematical correctness or independent source verification.
\end{itemize}
\end{document}
